\documentclass[11pt]{article}
\usepackage{series}
\usepackage{amsmath,amssymb,amsthm}
\usepackage{geometry}
\usepackage{hyperref}
\usepackage{enumitem}

\geometry{margin=1in}


% ---------- metadata ----------
\title{The Modal Triplet Theory Program B3:\\ Quantization as Discrete Constraint Encoding\\
in the Modal Triplet Theory Program}
\author{Peter Nero}
\date{January 2026}

\begin{document}

\maketitle

\begin{abstract}
We show that quantization arises as a necessary encoding class once reduced
descriptions encounter nil obstructions, i.e.\ termination of admissible
describability. In the Modal Triplet Theory (MTT) framework, quantization is not
introduced as a postulate, a correspondence principle, or a rule for promoting
classical observables to operators. Rather, it is identified as the unique
encoding response that preserves coherent structure when continuous descriptive
degrees of freedom are no longer admissible.

We demonstrate that nil obstructions force collapse of continuous families of
descriptions and select discrete survivors that are stable under admissible
refinement. These discrete structures constitute the content of quantization.
Topological and combinatorial invariants play a central role, while probability
appears only conditionally when invariant measures exist. No Hilbert space,
operator algebra, or measurement axiom is assumed at the structural level.

This paper completes the encoding-class triad of the Modal Triplet Theory
Program. Together with gravity as kinematic consistency encoding (circle) and
gauge structure as redundancy encoding (lens), quantization is shown to arise as
an inevitable response to termination of describability (nil).
\end{abstract}

\tableofcontents
\newpage

% ============================================================
\section{Introduction and Scope}

The Modal Triplet Theory Program establishes that reduced descriptions are
necessarily local, that global coherence is generically obstructed, and that
failure of global coherence admits an exhaustive classification into three
structural obstruction types: circle, lens, and nil. Previous papers in the
series analyzed the encoding responses forced by circle obstructions (gravity)
and lens obstructions (gauge structure).

The purpose of the present paper is to analyze the structural consequences of
\emph{nil} obstructions. A nil obstruction occurs when no admissible reduced
description exists on a region: the encoding fiber is empty and admissible
continuation terminates. Nil is neither inconsistency nor redundancy; it is the
failure of describability itself.

We show that nil obstructions force the introduction of a discrete constraint
encoding whose role is to preserve coherent structure when continuous
descriptions fail. This encoding is what is conventionally identified as
quantization. In the MTT framework, quantization is therefore not a modification
of classical dynamics, nor a statement about microscopic ontology. It is the
unique encoding response to termination of admissibility.

Throughout this paper we adopt the following principles:

\begin{itemize}
\item Quantization is an \emph{encoding}, not a dynamical rule.
\item Discreteness arises from structural constraints, not from axiomatic
postulates.
\item Probability is conditional and secondary; it appears only when invariant
measures exist.
\item Topological and combinatorial invariants play the primary role in selecting
admissible discrete structures.
\end{itemize}

The analysis is structural and applies to all realization classes compatible with
the MTT core. Familiar mathematical formalisms of quantum theory appear only as
realizations of the discrete constraint encoding and are not assumed in the
definition.

\begin{remark}[Position in the series]
This paper completes the encoding-class analysis forced by the
circle--lens--nil classification. It follows the gravity and gauge papers and
precedes works on unified encodings, concrete realizations, and phenomenological
interpretation. No additional structural assumptions beyond those of the MTT
core and obstruction classification are introduced.
\end{remark}

% ============================================================
% Next section will be:
% \section{Nil Obstructions and Collapse of Continuous Description}
% ============================================================
\section{Nil Obstructions and Collapse of Continuous Description}

In this section we analyze the structural consequences of nil obstructions for
reduced description. We show that nil obstructions force the collapse of
continuous descriptive degrees of freedom and necessitate a discrete constraint
encoding.

\subsection{Nil revisited}

Recall that a nil obstruction is defined as the absence of any admissible reduced
encoding on a region.

\begin{definition}[Nil obstruction]
A \emph{nil obstruction} exists at a point or region if the encoding fiber is
empty, i.e.\ no admissible local description applies.
\end{definition}

Nil is not inconsistency (circle) and not redundancy (lens); it is the failure of
describability itself.

\subsection{Termination of admissible continuation}

Nil obstructions manifest kinematically as termination.

\begin{lemma}
If a continuation chain encounters a nil obstruction, then no admissible
continuation exists beyond that point.
\end{lemma}

\begin{proof}
By definition of nil obstruction, no admissible encoding exists in that region.
Therefore no admissible continuation chain can be extended through or beyond it.
\end{proof}

This termination is structural and independent of any realization.

\subsection{Collapse of continuous families}

We now examine the effect of nil on continuous descriptive degrees of freedom.

\begin{theorem}[Collapse of continuous description]
In the presence of a nil obstruction, any continuous family of admissible reduced
descriptions collapses: only isolated or discrete descriptive structures may
remain admissible.
\end{theorem}

\begin{proof}
Assume a continuous family of admissible descriptions persists in a neighborhood
of a nil obstruction. Then admissibility would extend continuously into the nil
region, contradicting the definition of nil. Therefore continuity of description
cannot be maintained across a nil boundary. Only discrete survivors that do not
require continuity of admissibility can remain.
\end{proof}

\begin{remark}
This collapse does not require any notion of limit or divergence; it follows
purely from the incompatibility of continuity with absence of admissible
encoding.
\end{remark}

\subsection{Discrete survivors}

The collapse of continuous description selects discrete structures.

\begin{definition}[Discrete survivor]
A \emph{discrete survivor} is a reduced descriptive structure that remains
admissible under arbitrarily fine admissible refinement near a nil obstruction.
\end{definition}

Discrete survivors are isolated in the space of descriptions and are stable under
refinement.

\subsection{Stability under refinement}

Stability is the key distinguishing feature of discrete survivors.

\begin{lemma}
A reduced descriptive structure is a discrete survivor if and only if it is
stable under admissible refinement in the presence of a nil obstruction.
\end{lemma}

\begin{proof}
If a structure is stable under refinement, then arbitrarily fine refinement does
not introduce new admissible alternatives, so the structure persists as an
isolated description. Conversely, if a structure is not isolated, refinement
would generate nearby admissible descriptions, contradicting the collapse
imposed by nil.
\end{proof}

\subsection{No continuous resolution of nil}

One might attempt to resolve nil obstructions by modifying admissibility or
relaxing continuity.

\begin{lemma}
No modification of continuous descriptive structure can resolve a nil
obstruction without eliminating admissible description entirely.
\end{lemma}

\begin{proof}
Nil is defined as the absence of admissible description. Any attempt to extend
continuous description across a nil region either violates admissibility or
reintroduces an admissible encoding, contradicting the definition of nil.
Therefore continuous resolution is impossible.
\end{proof}

\subsection{Interpretation}

Nil obstructions therefore force a structural transition:

\begin{remark}
When nil obstructions are present, the theory must abandon continuous
descriptive degrees of freedom and restrict attention to discrete survivors.
This transition is not a dynamical process but a constraint on what descriptions
remain admissible.
\end{remark}

\subsection{Preview: quantization as encoding}

The necessity of discrete survivors motivates the definition of a discrete
constraint encoding.

\begin{remark}
The encoding that organizes and preserves discrete survivors under admissible
refinement is what is conventionally identified as quantization. In the next
section we define this encoding abstractly and show that it is uniquely forced by
nil obstructions.
\end{remark}

% ============================================================
% Next section will be:
% \section{Discrete Constraint Encoding and Quantization}
% ============================================================
\section{Discrete Constraint Encoding and Quantization}

We now define the encoding class forced by nil obstructions. This encoding
organizes the discrete survivors identified in the previous section and ensures
their stability under admissible refinement. In the Modal Triplet Theory
framework, this encoding is identified as quantization.

\subsection{Role of the discrete constraint encoding}

From the analysis of nil obstructions, the required encoding must satisfy the
following properties:

\begin{enumerate}[label=(\roman*)]
\item it must preserve discrete survivors under arbitrarily fine admissible
refinement;
\item it must exclude continuous degrees of freedom that are incompatible with
nil;
\item it must be definable locally on admissible domains and respect overlap
consistency;
\item it must not reduce to redundancy bookkeeping (lens) or kinematic
consistency bookkeeping (circle).
\end{enumerate}

Any encoding failing to satisfy these properties cannot resolve nil obstructions
without reintroducing inadmissible continuous structure.

\subsection{Definition of discrete constraint encoding}

We now define the encoding formally.

\begin{definition}[Discrete constraint encoding]
A \emph{discrete constraint encoding} is an admissible encoding class that
restricts reduced descriptions to discrete survivor sets selected by stability
under admissible refinement, and excludes all continuous descriptive degrees of
freedom incompatible with nil obstructions.
\end{definition}

The discrete constraint encoding does not generate discreteness; it records and
organizes the discreteness forced by nil.

\subsection{Resolution of nil obstructions}

We now show that discrete constraint encoding is necessary and sufficient to
resolve nil obstructions.

\begin{theorem}[Resolution of nil by discrete constraint encoding]
Nil obstructions can be resolved, in the sense of preserving coherent structure
across admissible domains, if and only if a discrete constraint encoding is
introduced.
\end{theorem}

\begin{proof}
(\emph{If}) Given a discrete constraint encoding, admissible description is
restricted to discrete survivors that remain stable under refinement near nil
boundaries. These survivors provide a consistent reduced description without
requiring continuity across nil regions.

(\emph{Only if}) In the absence of a discrete constraint encoding, any attempt to
preserve coherent structure across nil obstructions would require continuous
descriptive degrees of freedom, contradicting the collapse of continuity forced
by nil. Therefore no resolution is possible without such an encoding.
\end{proof}

\subsection{Uniqueness of the discrete constraint encoding}

We now establish the uniqueness of quantization as an encoding response to nil.

\begin{theorem}[Uniqueness of discrete constraint encoding]
Any admissible encoding that preserves coherent structure in the presence of nil
obstructions is equivalent, up to admissible re-encoding, to a discrete
constraint encoding.
\end{theorem}

\begin{proof}
Let $\mathcal Q$ be an admissible encoding resolving nil obstructions. By
definition, $\mathcal Q$ must restrict descriptions to structures stable under
arbitrarily fine admissible refinement. This restriction induces a discrete
selection on the space of possible descriptions. Any such encoding factors
through the quotient that identifies all inadmissible continuous variations,
which is precisely the discrete constraint encoding. Therefore $\mathcal Q$ is
equivalent to a discrete constraint encoding.
\end{proof}

\subsection{Quantization as encoding}

We are now in a position to identify quantization within the MTT framework.

\begin{remark}[Quantization as encoding]
The discrete constraint encoding defined above is what is conventionally
identified as quantization. In Modal Triplet Theory, quantization is not a rule
for promoting classical observables to operators, nor a statement about
microscopic ontology. It is the unique encoding response to termination of
describability imposed by nil obstructions.
\end{remark}

\subsection{No dynamical assumptions}

As with gravity and gauge encodings, quantization is defined structurally.

\begin{remark}
The discrete constraint encoding specifies \emph{what} descriptive structures
remain admissible in the presence of nil obstructions, not \emph{how} those
structures evolve. Dynamical laws and operator formalisms appear only in specific
realization classes and are not part of the encoding definition.
\end{remark}

\subsection{Preview: topology and combinatorics}

The selection of discrete survivors is closely tied to topological and
combinatorial invariants.

\begin{remark}
In the next section we show how topological and combinatorial structures
naturally organize discrete survivors and why such invariants play a central
role in quantized descriptions.
\end{remark}

% ============================================================
% Next section will be:
% \section{Topology and Combinatorial Invariants}
% ============================================================
\section{Topology and Combinatorial Invariants}

In this section we show that discrete survivors selected by nil obstructions are
naturally organized by topological and combinatorial invariants. These invariants
do not introduce new structure; they classify the discrete descriptive remnants
that remain admissible when continuous descriptions collapse.

\subsection{Why topology becomes relevant}

When continuous descriptive degrees of freedom are excluded by nil obstructions,
only structures invariant under admissible refinement can persist.

\begin{remark}
Refinement-invariant structures are necessarily discrete and global in nature.
Topology and combinatorics provide precisely such invariants: they remain well
defined under arbitrary local refinement and do not depend on metric or smooth
structure.
\end{remark}

Thus, topology enters quantization not as a geometric axiom, but as a consequence
of refinement stability.

\subsection{Discrete invariants under refinement}

We formalize the notion of refinement-invariant discrete data.

\begin{definition}[Refinement-invariant discrete invariant]
A \emph{refinement-invariant discrete invariant} is a discrete label or quantity
assigned to a reduced description such that:
\begin{enumerate}[label=(\roman*)]
\item it is preserved under admissible refinement of encodings;
\item it cannot be continuously deformed within the admissible description
space;
\item it distinguishes inequivalent discrete survivors.
\end{enumerate}
\end{definition}

Examples include winding numbers, intersection indices, combinatorial charges,
and other discrete classification data.

\subsection{Topological classification of survivors}

Discrete survivors are classified by their invariants.

\begin{theorem}[Topological organization of discrete survivors]
Discrete survivors selected by nil obstructions decompose into equivalence
classes labeled by refinement-invariant topological or combinatorial invariants.
\end{theorem}

\begin{proof}
Let $\mathcal S$ denote the set of discrete survivors near a nil boundary. Any two
elements of $\mathcal S$ that differ by a refinement-invariant discrete invariant
cannot be related by admissible re-encoding or refinement. Conversely, if no such
invariant distinguishes them, refinement stability implies they represent the
same survivor. Therefore $\mathcal S$ decomposes into equivalence classes labeled
by such invariants.
\end{proof}

This classification replaces continuous parameterization.

\subsection{Exclusion of continuous parameters}

Topology explains why continuous parameters are excluded.

\begin{lemma}
Any continuous parameter labeling discrete survivors is incompatible with
refinement stability in the presence of nil obstructions.
\end{lemma}

\begin{proof}
A continuous parameter would admit arbitrarily small variations. Under
refinement, such variations would generate nearby admissible descriptions,
contradicting the isolation required of discrete survivors. Therefore continuous
parameters cannot label admissible discrete survivors.
\end{proof}

\subsection{Relation to familiar quantum discreteness}

The discreteness encountered here encompasses, but does not assume, familiar
quantum phenomena.

\begin{remark}
Discrete spectra, quantized charges, and topological quantum numbers commonly
encountered in physical theories are realizations of refinement-invariant
discrete invariants. Their appearance is explained here as a structural
consequence of nil obstructions, not as a postulate of quantum theory.
\end{remark}

\subsection{Combinatorial structures}

Not all discrete invariants are topological in the geometric sense.

\begin{remark}
In some realizations, discrete survivors are organized by purely combinatorial
data (graph connectivity, adjacency relations, incidence structures) rather than
by geometric topology. The present framework accommodates both cases, as only
refinement invariance is required.
\end{remark}

\subsection{Independence from realization details}

The role of topology and combinatorics is structural, not realizational.

\begin{remark}
The appearance of topological or combinatorial invariants does not depend on the
existence of a smooth manifold or geometric background. Such invariants classify
discrete survivors in any realization class compatible with nil obstructions.
\end{remark}

\subsection{Preview: probability and measurement}

Discrete classification alone does not assign weights or probabilities to
survivors.

\begin{remark}
Probability arises only when additional structure—such as invariant measures on
discrete survivor sets—exists. In the absence of such structure, quantization
yields discrete possibilities without probabilistic interpretation. This issue is
addressed in the next section.
\end{remark}

% ============================================================
% Next section will be:
% \section{Conditional Probability and Measurement}
% ============================================================
\section{Conditional Probability and Measurement}

In this section we analyze the status of probability and measurement within the
discrete constraint encoding framework. We show that probability is not a
primitive ingredient of quantization, but a conditional structure that arises
only when additional invariant measures exist on the space of discrete
survivors.

\subsection{Absence of intrinsic probability}

Discrete survivors selected by nil obstructions are classified by refinement-
invariant discrete data. This classification alone does not assign probabilities.

\begin{remark}
Quantization, as defined by the discrete constraint encoding, yields a discrete
set of admissible descriptions but does not endow them with any intrinsic
probabilistic weighting. Discreteness does not imply randomness.
\end{remark}

Thus, probability is not forced by nil obstructions.

\subsection{Invariant measures on survivor sets}

Probability enters only when discrete survivor sets admit invariant measures.

\begin{definition}[Invariant measure on discrete survivors]
Let $\mathcal S$ be a discrete survivor set. An \emph{invariant measure} on
$\mathcal S$ is a probability measure $\mu$ such that $\mu$ is preserved under
all admissible re-encodings and refinement operations.
\end{definition}

Invariant measures are additional structures, not consequences of quantization.

\subsection{Conditional probability}

We now define probability conditionally.

\begin{definition}[Conditional probability]
Probability is defined on a discrete survivor set $\mathcal S$ if and only if an
invariant measure $\mu$ exists on $\mathcal S$. In that case, probabilities are
given by $\mu$; in the absence of such a measure, no probabilistic interpretation
is defined.
\end{definition}

\begin{remark}
This definition matches the treatment of probability in the structural core of
Modal Triplet Theory. Probability is conditional on the existence of invariant
measures and is not universally available.
\end{remark}

\subsection{Measurement as selection among survivors}

Measurement is not an axiom but a kinematic event.

\begin{definition}[Measurement]
A \emph{measurement} is a selection event in which admissible continuation
terminates and a particular discrete survivor is realized.
\end{definition}

Measurement corresponds to the encounter of a nil obstruction followed by
selection among discrete survivors.

\subsection{Measurement without collapse postulate}

No collapse postulate is assumed.

\begin{remark}
Selection among discrete survivors is forced structurally by admissibility
constraints. It does not require an additional collapse axiom or modification of
dynamics. The appearance of collapse is the manifestation of termination of
describability.
\end{remark}

\subsection{Repeatability and records}

Repeatable measurement outcomes correspond to stable discrete survivors.

\begin{lemma}
A discrete survivor that is stable under admissible refinement defines a
repeatable measurement outcome.
\end{lemma}

\begin{proof}
Stability under refinement ensures that subsequent admissible continuation cannot
alter the survivor label. Therefore the outcome persists and can be recorded.
\end{proof}

\subsection{Born-type rules as emergent}

When invariant measures exist, familiar probabilistic rules may emerge.

\begin{remark}
In specific realization classes, invariant measures may take forms that reproduce
Born-type probability rules. Such results are realizational and do not alter the
structural role of quantization as discrete constraint encoding.
\end{remark}

\subsection{No universal probability}

Finally, we emphasize the limits of probabilistic interpretation.

\begin{theorem}[No universal probability]
There exists no universally defined probability measure applicable to all
quantized descriptions in the Modal Triplet Theory framework.
\end{theorem}

\begin{proof}
Probability requires an invariant measure. Such measures may exist in some
discrete survivor sets but not in others. Therefore no universal probability
assignment is possible.
\end{proof}

\subsection{Preview: relation to other encodings}

Probability interacts with other encoding classes only indirectly.

\begin{remark}
Gauge and gravity encodings constrain which discrete survivors exist and how they
are organized, but they do not in themselves define probabilistic weights. The
interaction between probability, gauge, and gravity arises only in particular
realization classes.
\end{remark}

% ============================================================
% Next section will be:
% \section{Relation to Gravity and Gauge Encodings}
% ============================================================
\section{Relation to Gravity and Gauge Encodings}

In this section we clarify how the discrete constraint encoding associated with
quantization interacts with the kinematic consistency encoding (gravity) and the
redundancy encoding (gauge). Each encoding resolves a distinct obstruction type,
and their coexistence is structured rather than competitive.

\subsection{Distinct obstruction roles}

We recall the correspondence established in the Modal Triplet Theory framework:

\begin{center}
\begin{tabular}{c c c}
Obstruction & Encoding response & Structural role \\
\hline
Circle & Gravity & Kinematic consistency \\
Lens & Gauge & Redundancy bookkeeping \\
Nil & Quantization & Discrete constraint \\
\end{tabular}
\end{center}

Each encoding addresses a distinct failure of global coherence and therefore
cannot replace the others.

\subsection{Quantization does not resolve circle or lens}

Quantization resolves termination of describability, not inconsistency or
redundancy.

\begin{lemma}
The discrete constraint encoding does not resolve circle or lens obstructions.
\end{lemma}

\begin{proof}
Discrete constraint encoding restricts admissible descriptions to discrete
survivors. It does not modify overlap transport rules (required for circle) nor
does it identify redundant lifts (required for lens). Therefore it cannot restore
kinematic path independence or eliminate redundancy of representation.
\end{proof}

This establishes that quantization is structurally orthogonal to gravity and
gauge.

\subsection{Gravity and gauge constrain discrete survivors}

Although quantization does not resolve circle or lens, those encodings constrain
which discrete survivors are admissible.

\begin{remark}
Gravity encoding constrains discrete survivors by stabilizing causal structure
and kinematic continuation near nil boundaries. Gauge encoding constrains
discrete survivors by organizing redundancy of representation within each
survivor class.
\end{remark}

Thus, discrete constraint encoding operates within boundaries shaped by gravity
and gauge.

\subsection{Compatibility of encodings}

We now state the compatibility result.

\begin{theorem}[Compatibility of encoding classes]
The gravity, gauge, and quantization encodings are mutually compatible: they may
coexist without contradiction in a single realization, each addressing a
distinct obstruction type.
\end{theorem}

\begin{proof}
Each encoding modifies a different structural aspect of reduced description:
gravity modifies path-dependent transport, gauge modifies fiber redundancy, and
quantization modifies admissible descriptive sets. Since these aspects are
orthogonal, the encodings commute structurally and may be applied simultaneously.
\end{proof}

\subsection{Hierarchical interaction}

While compatible, the encodings interact hierarchically.

\begin{remark}
Gravity acts at the level of kinematic persistence and causal structure. Gauge
acts at the level of representational redundancy. Quantization acts at the level
of admissible descriptive existence. This hierarchy determines the order in which
constraints become active under refinement.
\end{remark}

This hierarchy explains why gravity and gauge often appear in classical regimes,
while quantization becomes dominant near nil boundaries.

\subsection{Mixed realizations}

Many familiar physical theories realize more than one encoding simultaneously.

\begin{remark}
Electromagnetism provides a realization in which gauge (lens) and circle
obstructions coexist within a single U(1) bundle, while quantization restricts
allowed charge and flux values discretely. Nonabelian gauge theories and
gravitational theories provide further examples of mixed realizations.
\end{remark}

The present framework accommodates such coexistence without conflation.

\subsection{No unification at the encoding level}

Finally, we emphasize the limits of unification.

\begin{remark}
Unification of gravity, gauge, and quantization does not occur at the level of
structural obstructions or encoding classes. It may occur only at the level of
specific realizations that simultaneously resolve multiple obstruction types.
Such unification is contingent and does not alter the distinct structural roles
identified here.
\end{remark}

\subsection{Preview: unified and saturated encodings}

The coexistence of multiple encodings motivates the study of saturated encoding
frameworks.

\begin{remark}
In subsequent papers we analyze encoding frameworks that simultaneously resolve
circle, lens, and nil obstructions in a single realization. Such frameworks
include, but are not limited to, string-theoretic and related constructions.
\end{remark}

% ============================================================
% Next section will be:
% \section{Summary and Outlook}
% ============================================================
\section{Summary and Outlook}

In this paper we have identified quantization as a necessary encoding class
arising from nil obstructions in the Modal Triplet Theory framework. Nil
obstructions correspond to termination of admissible describability: regions in
which no continuous reduced description can exist. Such obstructions are neither
inconsistency (circle) nor redundancy (lens), but failure of description itself.

We showed that nil obstructions force collapse of continuous descriptive degrees
of freedom and select discrete survivors that remain stable under admissible
refinement. The organization and preservation of these discrete survivors
requires a discrete constraint encoding, which we identified with quantization.
In this framework, quantization is not a postulate, a correspondence principle,
or a rule for operator promotion; it is the unique encoding response that allows
coherent structure to persist when continuity is structurally forbidden.

Topology and combinatorial invariants were shown to play a central role in
classifying discrete survivors. These invariants arise not from geometric
assumptions, but from refinement stability. Probability was shown to be secondary
and conditional: it appears only when invariant measures exist on discrete
survivor sets, and no universal probabilistic assignment is implied by
quantization itself. Measurement was interpreted as a selection event associated
with nil boundaries, without invoking collapse axioms or observer-dependent
postulates.

We further clarified the relation between quantization, gravity, and gauge
encodings. Each encoding resolves a distinct obstruction type—nil, circle, and
lens respectively—and the three are structurally compatible. Quantization does
not resolve kinematic path dependence or redundancy of representation, but
operates within the constraints imposed by gravity and gauge encodings. Their
interaction is hierarchical rather than competitive, explaining the coexistence
of classical, gauge, and quantum structures in familiar physical theories.

Together with the preceding papers on gravity as kinematic consistency encoding
and gauge structure as redundancy encoding, the present work completes the
encoding-class triad forced by the circle–lens–nil classification. These three
encodings exhaust the necessary structural responses to failure of global
coherence in reduced description.

Subsequent papers in the Modal Triplet Theory Program explore two directions.
First, unified or saturated encoding frameworks are analyzed, in which multiple
obstruction types are resolved simultaneously within a single realization.
Second, concrete realization papers construct explicit geometric, bundle-based,
and algebraic models that instantiate the encoding classes identified here,
without altering the structural conclusions of the present analysis.

In this way, quantization appears not as a mystery of microscopic physics, but as
an inevitable consequence of the limits of describability itself.

% ============================================================
% End of Paper B5

\end{document}
